\chapter{Context}
\label{ch:context}

\begin{figure}[H]
  \centering
  \resizebox{\textwidth}{!}{% The system context: Blendity, the people and programs around it.
\begin{tikzpicture}[node distance=9mm and 12mm]
  \node[box, fill=blAmber!15, draw=blAmber, minimum width=52mm, minimum height=22mm, align=center] (bl)
    {\textbf{Blendity}\\[1pt] \footnotesize one process, no plugins, no scripting runtime\\ \footnotesize the editor, \code{blendity\_tests}, \code{blendity\_stress}};
  \node[person, above=of bl] (artist) {The artist\\ \footnotesize keyboard, mouse, pen, fingers, a VR headset};
  \node[ext, left=18mm of bl, align=center] (os) {Windows / Linux / macOS\\ \footnotesize one window, input, files,\\ \footnotesize the trash, dialogs};
  \node[ext, below left=12mm and -6mm of bl, align=center] (gpu) {Any Vulkan GPU\\ \footnotesize (OptiX on NVIDIA)};
  \node[ext, below=of bl, align=center] (libs) {Blender's prebuilt libraries\\ \footnotesize optional, behind \code{BL\_WITH\_*}};
  \node[ext, below right=12mm and -6mm of bl, align=center] (dl) {Optional downloads\\ \footnotesize the libraries, AI inbetweening, MoonRay};
  \node[ext, right=18mm of bl, align=center] (engines) {Blender, Unity, Godot, Unreal\\ \footnotesize Live Link bridges (TCP, loopback)\\ \footnotesize files: OBJ, FBX, glTF, STL, PLY, USD,\\ \footnotesize .unitypackage, .tscn, .blend};
  \node[ext, above right=8mm and 2mm of bl, align=center] (gh) {GitHub\\ \footnotesize releases, the daily update check};
  \node[ext, above left=8mm and 2mm of bl, align=center] (proj) {The project folder\\ \footnotesize \code{Assets/}, scenes, materials, renders};
  \draw[flow, <->] (artist) -- (bl);
  \draw[flow, <->] (os) -- (bl);
  \draw[flow] (bl) -- (gpu);
  \draw[dflow] (libs) -- (bl);
  \draw[dflow] (dl) -- (bl);
  \draw[flow, <->] (bl) -- (engines);
  \draw[dflow] (gh) -- (bl);
  \draw[flow, <->] (proj) -- (bl);
\end{tikzpicture}
}
  \caption{The system context. Solid arrows are data both ways or out; dashed arrows are optional inputs Blendity runs without.}
  \label{fig:context}
\end{figure}

\section{The people}
One artist at a time, at a desk or on a 13-inch 2-in-1, with a keyboard and a mouse, a pen, fingers on the screen, a
pen display with a dial and express keys, or an OpenXR headset. The HCI work (\adr{0085}) names the artists it designs
for: people arriving from Unity, Unreal, Blender and Adobe's tools, keyboard-first and pen-first. The one developer is
the user, working with an AI (chapter~\ref{ch:sdlc}).

\section{The operating system}
\code{src/platform} is one window and its input on Win32, X11 (with XInput~2 for pens and touch) and Cocoa, plus the
clipboard, native file dialogs, the trash and DPI, mirroring Blender's GHOST layer \cite{Blender2026}. Everything above
it is portable C++17 built by \code{build.bat} on Windows (MSVC, found automatically) and \code{build.sh} on Linux and
macOS, with CMake optional (C-01). Pixels are \code{0xAARRGGBB} so all three can blit without conversion; the editor
redraws only on input, animation or timers and otherwise sleeps in the OS event wait.

\section{The GPU}
Any Vulkan 1.2 GPU: the loader is opened at run time (\code{VK\_NO\_PROTOTYPES}), so a machine without one runs on the
CPU. The GPU does three things, each with a CPU equivalent: the views are drawn by \code{gpu::Raster} with the CPU's
formulas in GLSL (\adr{0061}); the path tracer runs its kernel as a Vulkan compute shader (hardware ray queries where
they exist) or, on NVIDIA cards, through OptiX (\adr{0037}, \adr{0054}); and lighting bakes gather on it. CPU and GPUs
render together, claiming samples from one counter.

\section{Blender's libraries, optional}
When \code{blender/lib/<platform>} exists beside the project, the build links Blender's own prebuilt libraries
(Embree, OpenImageDenoise, oneTBB, Eigen, OpenSubdiv, OpenEXR, OpenColorIO, Manifold, meshoptimizer, Open PGL, Jolt,
Vulkan + shaderc, FFmpeg and others), each behind \code{BL\_WITH\_<NAME>} with a fallback (\adr{0001}, C-02).
\code{BLENDITY\_NO\_LIBS=1} forces the dependency-free build, which CI tests on every system; the installers ship both
programs and Blendity hands over to Blendity-full when the libraries load (\adr{0033}). The Console's \code{libs} prints
what a build has.

\section{Files and other programs}
A project is a folder with \code{Assets/} inside, as Unity's (\code{BlendityProject} in the home folder unless another is
opened). Scenes are text (\code{.scene}, Zstandard-compressed when the library is there), materials \code{.mat},
procedural textures \code{.ptex}, clips \code{.anim}, sequences \code{.seq}, layouts \code{.blendity-ui}. Models move
through OBJ, FBX, glTF, STL, PLY and USD, verified both ways against Blender 5.2 (R-06), and through Unity's
\code{.unitypackage}, \code{.prefab}, \code{.unity}, \code{.mat} and \code{.anim}, Godot's \code{.tscn} and Blender's
\code{.blend} (through the installed Blender). Live Link streams the scene to Blender, Unity, Godot and Unreal on the
same computer as it changes (chapter~\ref{ch:bridges}).

\section{The network}
Two uses, both optional and both the user's choice: a once-a-day check of GitHub's releases for a newer build
(\code{src/core/releases.cpp}; Preferences $>$ Updates), and Live Link's TCP server on 127.0.0.1 (no other computer can
connect; it is one way out). There is no telemetry and no cloud; the only learned model in the product is the optional,
local AI inbetweening download (\adr{0046}).

\section{The programs it follows}
The context has a second ring: the programs whose ideas the editor reproduces, each credited in chapter~\ref{ch:licensing}
and \code{CREDITS.md}. Unity for the editor's shape, Blender for the data and the tools, SketchUp, UModeler and
Plasticity for direct modeling, ZBrush and Sculptris for sculpting, Substance Painter for texturing, SketchBook, Photoshop
and Krita for drawing, Mixamo and the CATS plugin for rigging and avatars, Cascadeur for physics in animation, Premiere
for the sequencer, Lightroom for the Editing tab, Unity's Mecanim for retargeting. Where Blendity differs from them, the
user guide says so at the moment it matters.
